% Generated by tex/build_swebok_structure.py from tex/chapters/ch01/03_ソフトウェア要求の基礎.tex
\subsubsection{ソフトウェア製品要求とソフトウェアプロジェクト要求}

\term{ソフトウェア製品要求}は、完成するソフトウェアそのものが満たすべき性質である。画面、処理、データ、外部連携、性能、信頼性、セキュリティなどが含まれる。

\term{ソフトウェアプロジェクト要求}は、ソフトウェアを作るプロジェクトの進め方を制約する要求である。費用、納期、人員、テスト環境、データ移行、利用者教育、保守準備などが例である。プロジェクト要求は、プロジェクト憲章、契約、計画書などに記録されることが多い。

\par\smallskip\noindent\refstepcounter{table}\label{tab:ch01-03}\textbf{表 \thetable: 第01章 ソフトウェア要求 表03}\par\smallskip
\begin{tabularx}{0.97\linewidth}{@{}p{.28\linewidth}X@{}}
\toprule
区分 & 見分け方 \\
\midrule
\term{製品要求} & 完成したソフトウェアがどう振る舞うか、どの品質を満たすかを述べる。 \\
\term{プロジェクト要求} & 開発の費用、納期、人員、進め方、移行、教育などを縛る。 \\
\bottomrule
\end{tabularx}

\MiniBox{例}{「利用者は注文履歴を検索できる」は製品要求である。「本番移行は年度末までに完了する」はプロジェクト要求である。}
